\section{Đóng góp}

Đồ án này không nhằm đề xuất một bộ giải duy nhất thay thế toàn bộ các hướng tiếp cận hiện có, mà tập trung vào hai hướng bổ sung cho nhau trên cùng nền tảng đồ thị các tập lồi: tái triển khai và mở rộng formulation GCS-Bézier đã có, đồng thời đề xuất một formulation GCS-DMS mới để tích hợp lựa chọn vùng lồi với điều khiển tối ưu phi tuyến. Các đóng góp chính được tách thành hai nhóm như sau.

\textbf{Đóng góp 1: Tái triển khai và mở rộng thực nghiệm GCS-Bézier với nhiều chiến lược phân hoạch lồi.}
đồ án tái triển khai khuôn khổ GCS-Bézier cho hoạch định chuyển động quanh vật cản dựa trên Graph of Convex Sets và đường cong Bézier~\cite{MarcucciEtAl2023GCS,MarcucciEtAl2023MotionPlanning}. Phần triển khai này bao gồm xây dựng đồ thị vùng lồi, mã hóa ràng buộc an toàn thông qua tính chất bao lồi của đường cong Bézier, áp dụng nới lỏng phối cảnh để đưa bài toán về dạng SOCP/MISOCP, và đánh giá nghiệm trên các môi trường 2D và mê cung.

Điểm mở rộng chính so với việc chỉ tái hiện lại công thức của Marcucci và cộng sự nằm ở phần phân hoạch không gian tự do. Thay vì chỉ dùng một cơ chế sinh vùng lồi cố định, đồ án khảo sát và so sánh ba họ thuật toán phân hoạch: ACD~\cite{LienAmato2006}, IRIS~\cite{DeitsTedrake2014} và VCC~\cite{DeitsTedrake2023VCC}. Trong đó, việc đưa ACD vào pipeline GCS-Bézier là một đóng góp mang tính engineering đáng kể: ACD khai thác trực tiếp hình học đa giác 2D, thường tạo số vùng ít hơn và ổn định hơn trong các môi trường phẳng, trong khi hướng GCS-Bézier nguyên bản của Marcucci chủ yếu gắn với các vùng lồi sinh bởi IRIS hoặc các biến thể liên quan. Phần thực nghiệm vì vậy không chỉ đánh giá bộ giải GCS-Bézier, mà còn làm rõ ảnh hưởng của chất lượng phân hoạch lên số vùng, số cạnh đồ thị, thời gian giải và chất lượng quỹ đạo.

\textbf{Đóng góp 2: Đề xuất formulation GCS-DMS với ràng buộc động lực học phi tuyến và hàm cản log-barrier.}
Đóng góp mới hơn của đồ án là mô hình GCS-DMS, trong đó bài toán hoạch định được viết dưới dạng điều khiển tối ưu trên đồ thị các vùng lồi. Khác với GCS-Bézier, quỹ đạo không được tham số hóa bằng các điểm điều khiển Bézier thuần hình học, mà bằng các biến trạng thái, điều khiển và thời lượng cục bộ trên từng vùng. Phương trình động lực học phi tuyến được rời rạc hóa bằng Direct Multiple Shooting, với các ràng buộc defect đảm bảo tính nhất quán giữa nghiệm tích phân cục bộ và trạng thái biên của từng vùng. Trong đó, biến chọn cạnh/vùng của GCS, biến trạng thái biên, biến điều khiển, thời lượng đoạn, ràng buộc ghép nối giữa các vùng và ràng buộc defect của multiple shooting được đặt trong một formulation chung thuộc lớp MIOCP/MINLP. Nhờ đó, bài toán vừa giữ được cấu trúc chọn đường đi rời rạc của GCS, vừa mô hình hóa trực tiếp động lực học của hệ robot.

Ngoài ra, đồ án đưa hàm cản logarit vào formulation GCS-DMS như một cơ chế đảm bảo an toàn hình học theo miền nội của từng polytope. Thành phần log-barrier phạt mạnh khi quỹ đạo tiến sát biên vùng an toàn, giúp nghiệm có xu hướng nằm sâu trong vùng khả thi thay vì chỉ thỏa ràng buộc tại một số điểm lưới hữu hạn. Đây là điểm khác biệt quan trọng so với các cách tích hợp DMS với GCS theo kiểu hậu xử lý hoặc chỉ áp đặt containment bằng sampling rời rạc. Formulation này đồng thời làm rõ sự đánh đổi: ta mất cấu trúc lồi MISOCP và bảo đảm tối ưu toàn cục của GCS-Bézier, nhưng đổi lại thu được một mô hình có khả năng kiểm chứng trực tiếp tính khả thi động lực học thông qua defect dynamics, connection gap, integrality gap và mức vi phạm ràng buộc.
